\chapter{Introduction}
\label{cha:introduction}

\section{Context and Motivation}
\label{sec:introduction:context}

A failure in a distributed application can become visible in a component
other than the one in which it originates. Research on microservice
performance debugging shows how problems propagate through dependencies,
complicating the identification of their source \parencite{gan2021sage}.
Diagnosing such failures therefore requires relating observations across
components and understanding how their interactions contribute to the
observed behaviour.

This challenge provides the operational context for the present thesis,
conducted in collaboration with Norlys. The company's \emph{Operations with
AI Agents} specification describes an application request path through
Cloudflare, Azure Firewall, Azure infrastructure and the application itself.
According to the specification, investigating a failing endpoint involves
consulting these systems separately and manually correlating timestamps and
request identifiers \parencite[Background]{norlys_operations_agents}.
The task is to assemble this distributed information into an explanation of
the failure that can support an appropriate corrective action.

Correlation provides a basis for connecting these observations. The W3C Trace
Context standard, for example, defines identifiers that can be propagated
between participating systems to associate their activity with a common
distributed trace \parencite{w3c2021tracecontext}. The interpretation of the
resulting evidence also depends on which information is relevant to the
problem being investigated. Research on machine learning in security shows that models can rely on patterns unrelated to the underlying security problem, leading to misleading assessments of their capabilities.
\parencite{arp2022dosdonts}. For the present project, this motivates examining
both the available telemetry and the features derived from it as potential
evidence for security-related diagnosis.

Large language model (LLM) agents offer a possible approach to coordinating
the investigation across these sources. ReAct combines model reasoning with
actions that retrieve observations from external environments
\parencite{yao2023react}, while AutoGen supports interaction between agents
with different roles, tools and levels of human participation
\parencite{wu2023autogen}. These approaches motivate investigating whether
specialised agents can gather and interpret telemetry from different layers
of the request path and combine their findings into a supported diagnosis.
The Norlys specification requires conclusions to be accompanied by evidence
and allows an inconclusive outcome when the available information is
insufficient \parencite[AC-4]{norlys_operations_agents}.

Extending this process to remediation introduces a further requirement:
changes must be controlled and their effects assessed. The NIST AI Risk
Management Framework emphasises testing in the intended context, monitoring
and human intervention \parencite[pp.~14--15]{tabassi2023airmf}. For this
project, Norlys specifies an action allowlist, dry-run validation, human
approval, recovery verification and rollback
\parencite[AC-5]{norlys_operations_agents}. Whether these controls
enable effective remediation while preventing unintended changes forms part
of the investigation.

The objective of this thesis is to design and evaluate a multi-agent proof
of concept using real operational telemetry from the selected
Cloudflare-to-Azure request path. The investigation brings together telemetry
correlation, the identification of security-relevant data and derived
features, evidence-based diagnosis and human-approved remediation. A
systematic literature review will provide the research basis for the work.

\section{Problem Statement and Research Questions}
\label{sec:introduction:research-questions}

This thesis addresses the following problem statement:

\begin{quote}
How can a multi-agent AI system be designed and evaluated to correlate
telemetry from multiple operational systems and support evidence-based
diagnosis and safe, human-approved remediation of controlled failures in a
non-production Cloudflare-to-Azure application request path?
\end{quote}

The investigation is structured around four research questions:

\begin{description}
    \item[RQ1:] How can a multi-agent system access and correlate telemetry
    from Cloudflare, Azure Firewall, Azure infrastructure, and applications
    to reconstruct the path of a failing endpoint request?

    \item[RQ2:] Which telemetry data and derived features are relevant for
    diagnosing security-related problems along the Cloudflare-to-Azure
    application request path?

    \item[RQ3:] How accurately and consistently can the system identify the
    faulty layer/component and support its diagnosis with relevant evidence
    across different scenarios?

    \item[RQ4:] How safely and effectively can the system plan and execute
    allow-listed remediations under human approval, including dry-run
    validation, recovery verification and rollback?
\end{description}

\section{Scope and Delimitations}
\label{sec:introduction:scope}

Norlys' wider \emph{Centralized Observability Platform Story} specifies a
platform for telemetry ingestion, normalisation, protection of sensitive data,
analysis and remediation. It explicitly states:
\enquote{This is a large cross-cutting story and should be decomposed into
smaller tasks and delivery milestones.}
\parencite[Description]{norlys_observability_spec}
The thesis addresses the focused \emph{Operations with AI Agents} proof of
concept within this wider initiative. Its scope covers investigation across
the four specified telemetry sources and human-approved remediation within
a designated test path. Platform-wide ingestion, normalisation, data
obfuscation and retention remain responsibilities of the parent initiative
\parencite[Scope]{norlys_operations_agents}.

The security-related investigation concerns telemetry and derived features
associated with the same request path. The findings will be interpreted in
relation to the available data, the selected environment and the scenarios
examined. These boundaries also frame the assessment of diagnostic
effectiveness and remediation safety.

\section{Main Contributions}
\label{sec:introduction:contributions}

% TODO: Complete after implementation and evaluation.

\section{Outline of the Thesis}
\label{sec:introduction:outline}

% TODO: Complete when the chapter structure is stable.